% ------------------------------------------------------------------------
%
\begin{frame}{Tree traversals}

  Let~$t$ be the tree $\fun{Int}(x, left, right)$.
  \begin{itemize}

    \item The \textbf{preorder traversal} of~$t$ is the list whose
      first element is~$x$, followed by the elements of the preorder
      traversal of~\emph{left}, followed by the elements of the
      preorder traversal of~\emph{right}. The previous tree yields
      $[8; 1; 3; 5; 9; 3; 2]$.

    \item The \textbf{postorder traversal} of~$t$ is the list made of
      the elements of the postorder traversal of~\emph{left}, followed
      by the elements of the postorder traversal of~\emph{right},
      followed by~$x$. The previous tree yields $[5; 3; 9; 1; 2; 3;
        8]$.

    \item The \textbf{inorder traversal} of~$t$ is the list made of
      the elements of the inorder traversal of~\emph{left}, followed
      by~$x$, followed by the elements of the inorder traversal
      of~\emph{right}. The previous tree yields $[3; 5; 1; 9; 8; 3;
        2]$.

  \end{itemize}

\end{frame}

% ------------------------------------------------------------------------
%
\begin{frame}{Tree traversals}

  We want to build the traversals only by pushing values on a stack.
  We must add an \textbf{auxiliary parameter}, originally set to the
  empty list, on which the contents of the nodes are pushed in the
  proper order. For preorder:
  \begin{equation*}
    \begin{array}{r@{\;}l@{\;}l@{}}
      \fun{pre}(t,s) & \rightarrow & \ldots\\
      \fun{pre'} \; t & \rightarrow & \fun{pre}(t, \el).
    \end{array}
  \end{equation*}
  This extra parameter is the initial value of an accumulator.  We
  should interpret this accumulator when considering the call
  \(\fun{pre} (t, s)\). Let us consider the tree \(t = \fun{Int} (x,
  t_1, t_2)\) in the following figures:
  \begin{center}
    \includegraphics[bb=71 672 165 722]{pre_tree}
    \includegraphics[bb=71 672 180 722]{in_tree}
    \includegraphics[bb=71 672 180 722]{post_tree}
  \end{center}
  \centerline{Preorder (\fun{pre}) \qquad\qquad Inorder
    (\fun{in}) \qquad Postorder (\fun{post})}

\end{frame}

% ------------------------------------------------------------------------
%
\begin{frame}{Preorder traversal}

  Each arrow is a step in the traversal and connects different stages
  of the accumulator: a downwards arrow points to the argument of a
  recursive call on the corresponding child; an upwards arrow points
  to the value of the call on the parent. The rules for external
  nodes are not shown and simply consist in leaving the accumulator
  invariant.

  \bigskip

  In the preorder traversal figure, the subtree~\(t_2\) is used in the
  first recursive call \(\fun{pre}(t_2, s)\) whose value is
  named~\(s_1\): \(\fun{pre}(t_2, s) \twoheadrightarrow s_1\).
  Likewise, we have \(\fun{pre}(t_1, s_1) \twoheadrightarrow
  s_2\). Therefore
  \begin{equation*}
    \fun{pre}(t_1, \fun{pre}(t_2, s)) \twoheadrightarrow s_2.
  \end{equation*}

  The root is associated with the evaluation
  $\fun{pre}(t, s) \twoheadrightarrow \cons{x}{s_2}.$
  Consequently,
  \begin{equation*}
    \fun{pre}(t, s) \equiv \cons{x}{\fun{pre}(t_1, \fun{pre}(t_2, s))}.
  \end{equation*}

\end{frame}

% ------------------------------------------------------------------------
%
\begin{frame}{Three traversals}

  Finally:
  \begin{equation*}
    \begin{array}{r@{\;}l@{\;}l@{}}
      \fun{pre}(\fun{Ext}(), s) & \rightarrow & s; \\
      \fun{pre}(\fun{Int}(x, t_1, t_2), s) & \rightarrow
      & \cons{x}{\fun{pre}(t_1, \fun{pre}(t_2, s))}.\\
      \fun{pre'} \; t & \rightarrow & \fun{pre}(t, \el).
    \end{array}
  \end{equation*}

  Similarly:
  \begin{equation*}
    \begin{array}{r@{\;}l@{\;}l@{}}
      \fun{in}(\fun{Ext}(), s) & \rightarrow & s; \\
      \fun{in}(\fun{Int}(x, t_1, t_2), s) & \rightarrow
      & \fun{in}(t_1, \cons{x}{\fun{in}(t_2, s)}).\\
      \fun{in'} \; t & \rightarrow & \fun{in}(t, \el).
    \end{array}
  \end{equation*}
  And
  \begin{equation*}
    \begin{array}{r@{\;}l@{\;}l@{}}
      \fun{post}(\fun{Ext}(), s) & \rightarrow & s; \\
      \fun{post}(\fun{Int}(x, t_1, t_2), s) & \rightarrow
      & \fun{post}(t_1, \fun{post}(t_2, \cons{x}{s})).\\
      \fun{post'} \; t & \rightarrow & \fun{post}(t, \el).
    \end{array}
  \end{equation*}

\end{frame}

% ------------------------------------------------------------------------
%
\begin{frame}{Three folds on binary trees}

  We can abstract the pushing and parameterise pre-, in- and
  post\hyp{}order traversals, so they act like a fold on trees:
  \begin{equation*}
    \begin{array}{r@{\;}l@{\;}l@{}}
      \fun{fold\_pre}(f, \fun{Ext}(), s) & \rightarrow & s; \\
      \fun{fold\_pre}(f, \fun{Int}(x, t_1, t_2), s) & \rightarrow
      & f(x, \fun{fold\_pre}(f, t_1, \fun{fold\_pre}(f, t_2, s)))).\\
      \fun{fold\_pre'}(f, t) & \rightarrow & \fun{fold\_pre}(f, t, \el).\\
      \fun{pre} \; t & \rightarrow & \fun{fold\_pre'}((\Xfun \; (x, s) \rightarrow \cons{x}{s}), t).\\\\
      \fun{fold\_in}(f, \fun{Ext}(), s) & \rightarrow & s; \\
      \fun{fold\_in}(f, \fun{Int}(x, t_1, t_2), s) & \rightarrow
      & \fun{fold\_in}(f, t_1, f (x, \fun{fold\_in}(f, t_2, s))).\\
      \fun{fold\_in'}(f, t) & \rightarrow & \fun{fold\_in}(f, t, \el).\\
      \fun{in} \; t & \rightarrow & \fun{fold\_in'}((\Xfun \; (x, s) \rightarrow \cons{x}{s}), t).\\\\
      \fun{fold\_post}(f, \fun{Ext}(), s) & \rightarrow & s; \\
      \fun{fold\_post}(f, \fun{Int}(x, t_1, t_2), s) & \rightarrow
      & \fun{fold\_post}(f, t_1, \fun{fold\_post}(f, t_2, f(x,s))).\\
      \fun{fold\_post'} (f, t) & \rightarrow & \fun{fold\_post}(f, t, \el).\\
      \fun{post} \; t & \rightarrow & \fun{fold\_post'}((\Xfun \; (x, s) \rightarrow \cons{x}{s}), t).
    \end{array}
  \end{equation*}

\end{frame}
